% !TeX program = lualatex
% Compile twice: lualatex open_problems_500.tex
% All references and prose are embedded; no BibTeX database or external inputs.
\documentclass[11pt,a4paper]{article}
\usepackage[margin=24mm,headheight=14pt]{geometry}
\usepackage{fontspec}
\setmainfont{Latin Modern Roman}
\setsansfont{Latin Modern Sans}
\setmonofont{Latin Modern Mono}
\usepackage{amsmath,amssymb,mathtools}
\usepackage{unicode-math}
\setmathfont{Latin Modern Math}
\usepackage{microtype}
\usepackage{enumitem,xurl,needspace}
\usepackage[dvipsnames]{xcolor}
\usepackage{hyperref}
\hypersetup{unicode=true,colorlinks=true,linkcolor=MidnightBlue,urlcolor=MidnightBlue,
 pdftitle={500 Mathematical Problem Statements},pdfauthor={Source-annotated compilation},bookmarksnumbered=true}
\usepackage{bookmark}
\usepackage{tocloft}
\setlength{\cftsecnumwidth}{3em}
\cftsetpnumwidth{2.5em}
\cftsetrmarg{4em}
\usepackage{titlesec}
\titleformat{\section}{\Large\bfseries\raggedright}{\thesection}{0.7em}{}
\usepackage{fancyhdr}
\pagestyle{fancy}\fancyhf{}
\fancyhead[L]{\small\sffamily Mathematical problem statements}
\fancyhead[R]{\small Source edition 22}
\fancyfoot[C]{\thepage}
\setlength{\parindent}{0pt}
\setlength{\parskip}{5pt plus 1pt}
\setlength{\emergencystretch}{3em}
\setcounter{tocdepth}{1}
\setcounter{secnumdepth}{1}
\setlist[itemize]{leftmargin=3em,itemsep=5pt,topsep=4pt,parsep=0pt}
\allowdisplaybreaks
\newcommand{\N}{\mathbb{N}}
\newcommand{\Z}{\mathbb{Z}}
\newcommand{\Q}{\mathbb{Q}}
\newcommand{\R}{\mathbb{R}}
\newcommand{\C}{\mathbb{C}}
\newcommand{\F}{\mathbb{F}}
\newcommand{\A}{\mathbb{A}}
\newcommand{\PP}{\mathbb P}
\newcommand{\E}{\mathbb{E}}
\newcommand{\eps}{\varepsilon}
\newcommand{\dd}{\,\mathrm d}
\newcommand{\sref}[2]{\hyperlink{source:#1:#2}{[S#2]}}
\newcommand{\eref}[2]{\hyperlink{extra:#1:#2}{[E#2]}}
% Turn the next switch off for a shorter reading copy. The quotations preserve
% every exactTarget field, including qualifications omitted from a short summary.
\newif\ifshowcatalogue
\showcataloguetrue
\newcommand{\cataloguescope}[1]{\ifshowcatalogue
 \paragraph{Exact scope in the supplied catalogue.}
 \begin{quote}\small #1\end{quote}\fi}

% Standalone research notebook, original catalogue problem 191.
% Compile twice with: lualatex 191.tex
\numberwithin{equation}{section}
\hypersetup{pdftitle={Research attempt -- problem 191},pdfauthor={Research notebook; original compilation retained}}
\fancyhead[L]{\small\sffamily Research notebook}
\fancyhead[R]{\small\sffamily Problem 191 | 27 September 2026}
\begin{document}
\setcounter{section}{190}
% ================================================================
% problemId: problem.kimura-osullivan-finite-dimensionality-conjecture
% source releaseRank: 191; source releaseStatus: open_with_solved_subcases
% exactTarget SHA-256: 6ee9c8ec43e58498e451dd53a6637a35ab28436f112ed90be2b801d7f3186328
\clearpage
\section{\texorpdfstring{Kimura–O'Sullivan Finite-Dimensionality Conjecture}{Kimura–O'Sullivan Finite-Dimensionality Conjecture}}\label{191}
\subsection{Definitions and mathematical statement}
In the rational Chow-motive category $\mathrm{CHM}(k)_{\Q}$, morphisms are rational correspondences modulo rational equivalence, idempotents are split, and tensor product uses products of varieties. Exterior and symmetric powers are the images of the alternating and averaging idempotents of symmetric groups under the ordinary symmetry constraint. The conjecture asserts
\[
 h(X)\cong M_+\oplus M_-,\qquad
 \Lambda^n M_+=0,\qquad\operatorname{Sym}^m M_-=0
\]
for every smooth projective $X/k$, for some $n,m\ge1$. Either summand may be zero. This categorical finite-dimensionality is not simply finiteness of a cohomology vector space. The coefficients are rational but the base field has arbitrary characteristic; no uniform exponents or effective splitting are requested.
\par\smallskip\textit{Formulation and concept references:} \sref{191}{1}.
\cataloguescope{For every field k and every smooth projective variety X over k, the Chow motive h(X) with rational coefficients in the ordinary symmetric, pseudo-abelian category CHM(k)\_Q is Kimura–O'Sullivan finite-dimensional: there exist motives M+ and M\ensuremath{-} and positive integers n,m such that h(X) is isomorphic to M+ direct-sum M\ensuremath{-}, the exterior power \ensuremath{\Lambda}\textasciicircum{}n(M+) is zero, and the symmetric power Sym\textasciicircum{}m(M\ensuremath{-}) is zero. The summands may be zero. The field k is unrestricted; characteristic zero applies to the rational coefficient field, not to k. No effective construction, uniform n,m or classification is required.}
\subsection{Short English statement}
Can every smooth projective variety's rational motive be split into an even part with a vanishing exterior power and an odd part with a vanishing symmetric power?
% BEGIN RESEARCH ADDITION ATTEMPT-191
\subsection{Research attempt}

\paragraph{Current frontier.}
Finite-dimensionality is known for motives generated from curves by the usual tensor operations and summands, but not for arbitrary smooth projective varieties. Its nilpotence consequences are central in the foundational account \sref{191}{1}. Recent results still concern special families, for example a nine-dimensional family of K3 surfaces \eref{191}{1}. We do not infer categorical finite-dimensionality from finite Betti numbers. The proposed route is to lift the cohomological parity decomposition and then kill the first unrealized Schur projectors; this will expose exactly how much nilpotence is needed.

\paragraph{Attack route.}
One route seeks explicit algebraic even/odd projectors; another tries to rule out homologically invisible direct summands of tensor powers. Combining them is promising because the Schur projector is already idempotent: any nilpotence theorem applicable to it forces actual vanishing. But this requires nilpotence on tensor powers, not merely on $h(X)$ itself. We give the full conditional construction, including a closed finite formula for lifting an approximate idempotent.

\paragraph{Attempt.}
Fix $X/k$ and a Weil realization $H$ with characteristic-zero coefficient field; in characteristic $p>0$ one may use an $\ell$-adic realization with $\ell\ne p$. Cohomology is treated as a super vector space with parity given by cohomological degree. Put $M=h(X)$ and $b_\pm=\dim H^{\rm even/odd}(M)$.

Assume the following two inputs for this $X$:
\begin{enumerate}[label=(\roman*),leftmargin=2.5em]
\item There is $q\in\operatorname{End}(M)$ whose realization is the even cohomological projector.
\item Every homologically trivial endomorphism of $M^{\otimes r}$ is nilpotent under composition, for each positive integer $r$ used below.
\end{enumerate}
These are hypotheses, not results of this notebook. It suffices in (ii) to use $r=1,b_++1,b_-+1$ (repeated or vacuous cases may be omitted).

Let $d=q^2-q$. By (i), $H(d)=0$, so (ii) gives $d^N=0$ for some $N\ge1$. All polynomials in $q$ commute. Define
\begin{equation}\label{eq191:lift}
 v=\sum_{j=0}^{N-1}\binom{-1/2}{j}(4d)^j,
 \qquad p=\frac{1+(2q-1)v}{2}.
\end{equation}
The formal power-series identity $(1+4d)^{-1/2}$ becomes a finite identity modulo $d^N$; hence
\[
 v^2(1+4d)=1,\qquad ((2q-1)v)^2=1,
 \qquad p^2=p.
\]
Moreover $H(v)=1$ and $H(p)=H(q)$. Thus the pseudo-abelian category splits
\[
 M=M_+\oplus M_-,\qquad M_+=(M,p),\quad M_-=(M,1-p),
\]
with the intended even and odd realizations. No infinite series converges in a Chow group: the nilpotence exponent made the formula finite.

Now set $r=b_++1$. On $M^{\otimes r}$ consider
\[
 e_+=p^{\otimes r}\left(\frac1{r!}\sum_{\sigma\in S_r}
             \operatorname{sgn}(\sigma)\sigma\right).
\]
The factors commute, so $e_+$ is idempotent. Its image is $\Lambda^rM_+$. On even realization spaces the symmetry has no Koszul sign, so $H(e_+)$ projects onto the ordinary exterior power of a $b_+$-dimensional space and is zero. Hypothesis (ii) gives $e_+^a=0$ for some $a$, while idempotence gives $e_+^a=e_+$. Hence $e_+=0$ and $\Lambda^{b_++1}M_+=0$.

For $r=b_-+1$ the relevant idempotent is
\[
 e_-=(1-p)^{\otimes r}\left(\frac1{r!}\sum_{\sigma\in S_r}\sigma\right).
\]
On an odd super vector space a permutation acts with its sign. Consequently the \emph{ordinary categorical symmetric projector} realizes as the ordinary exterior projector on the underlying odd vector space. Its realization is zero, so the same nilpotence/idempotence argument proves $\operatorname{Sym}^{b_-+1}M_-=0$. This sign check is essential; replacing the motive category's symmetry by an unsigned symmetry on all cohomology would make the argument false. If $b_+=0$ or $b_-=0$, the corresponding exponent is one and the same reasoning forces that summand itself to vanish.

We have therefore proved, under (i)--(ii), precisely the finite-dimensionality decomposition requested in the problem, with exponents determined by the chosen realization. Rational coefficients permit the denominators $2$, $r!$ and the binomial coefficients; they impose no restriction on the characteristic of $k$.

\paragraph{Stress test.}
The first hypothesis is an algebraicity statement for the parity projector, not a formal consequence of K\"unneth decomposition in cohomology. The second is a tensor-power nilpotence assertion. Knowing nilpotence only for $\operatorname{End}(M)$ permits the lift \eqref{eq191:lift} but does not justify killing $e_\pm$, which live on larger powers. A nonzero homologically trivial idempotent on such a power is exactly the phantom summand this route must exclude.

Conversely, one cannot cite Kimura's nilpotence theorem to establish (ii) for the unknown motive and then conclude finite-dimensionality: that would use finite-dimensionality as a premise of its own proof. The known implications and the distinction between numerical, homological and smash nilpotence are discussed in \sref{191}{1}; no interchange of composition and external products is made here. The symbolic checks verify the truncated inverse-square-root identity modulo $d^N$ for small $N$, while the formal identity above proves it for every $N$. Neither check provides either missing geometric hypothesis.

\paragraph{Outcome.}
\textbf{Outcome: Conditional reduction.}
Algebraicity of the even projector plus homological nilpotence on the specified tensor powers yields the complete requested splitting by an explicit finite construction. The proof identifies two independent unresolved inputs and the exact point where tensor powers matter. No new unconditional finite-dimensionality class, nor a claim of priority for the general lifting principle, is asserted.

% END RESEARCH ADDITION ATTEMPT-191
\subsection{Sources}
\begin{itemize}\raggedright
\item[\textnormal{[S1]}]\hypertarget{source:191:1}{} Motifs de dimension finie [d'après S.-I. Kimura, P. O'Sullivan, ...]; Bourbaki929,March2004; Astérisque299(2005),115–145.
\url{https://www.numdam.org/item/SB_2003-2004__46__115_0.pdf}.
{\small\textit{Catalogue formulation source.}}
% BEGIN RESEARCH ADDITION REFERENCES-191
\item[\textnormal{[E1]}]\hypertarget{extra:191:1}{} Michele Bolognesi and Robert Laterveer, \emph{A 9-dimensional family of K3 surfaces with finite-dimensional motive}, Rendiconti del Circolo Matematico di Palermo, Series 2, 73 (2024), 2313--2331; arXiv:2310.11981v2. Theorem 3.1 and the spread argument in Section 2.5.
\url{https://arxiv.org/html/2310.11981v2}.
{\small\textit{Research reference; consulted 27 September 2026.}}
% END RESEARCH ADDITION REFERENCES-191
\end{itemize}

\end{document}
